\section{Badane protokoły}

\subsection{Kryteria doboru}

\subsubsection{Kryterium nadrzędne}
Rozwiązanie musi stanowić samoistną, bezpośrednią implementację mechanizmów tunelowania i kryptografii. Wyklucza się nakładki oprogramowania oraz systemy orkiestracji.

\subsubsection{Kryteria opisowe}

\begin{enumerate}[leftmargin=*]
    \item \textbf{Dojrzałość technologiczna i stabilność} - protokół powinien posiadać stabilną wersję oraz udokumentowaną specyfikację.
    \item \textbf{Dostępność i otwartość kodu źródłowego} - preferowane są rozwiązania otwartoźródłowe.
    \item \textbf{Powszechność zastosowania} - protokół powinien mieć rzeczywiste wdrożenia produkcyjne lub szeroką bazę użytkowników.
    \item \textbf{Aktualność i perspektywa rozwoju} - protokoł nie jest wycofany z użycia lub uznany za przestarzały.
    \item \textbf{Wsparcie dla różnych platform} - protokół powinien posiadać wsparcie dla różnych architektur i systemów operacyjnych.
\end{enumerate}

\subsubsection{Kryteria kwalifikacyjne}
\begin{enumerate}[leftmargin=*]
    \item Z pominięciem oprogramowania klienckiego w modelu scentralizowanym, wszelkie oprogramowanie wymagane do wdrożenia infrastruktury danego protokołu udostępnione jest na licencji zatwierdzonej przez OSI. \cite{53EK}
    \item Posiada opublikowaną specyfikację techniczną.
    \item Posiada co najmniej jedną stabilną wersję oznaczoną jako produkcyjna.
    \item Jest aktywnie rozwijany, co potwierdza aktywność w repozytorium kodu źródłowego w ciągu ostatnich 24 miesięcy.
    \item Oprogramowanie serwerowe protokołu posiada wsparcie dla systemu operacyjnego Linux.
    \item Wdrożenie i utrzymanie nie wymaga użycia infrastruktury lub usług stron trzecich.
\end{enumerate}

\newpage

\subsubsection{Wyniki kwalifikacji}

Na podstawie zdefiniowanych kryteriów kwalifikacyjnych, do dalszej analizy zakwalifikowano następujące rozwiązania:

\begin{enumerate}[leftmargin=*]
    \item \textbf{OpenVPN} - dojrzałe rozwiązanie realizujące model scentralizowany. Stanowi powszechnie stosowany punkt odniesienia w~pracach porównawczych dotyczących sieci VPN.
    \item \textbf{WireGuard} - minimalistyczny protokół warstwy trzeciej. W~zależności od konfiguracji realizuje model scentralizowany lub rozproszony.
    \item \textbf{Nebula} - protokół mesh VPN oparty na modelu rozproszonym z~węzłami lighthouse pełniącymi rolę koordynacyjną. Całość infrastruktury utrzymywana jest przez użytkownika.
\end{enumerate}

\newpage
\subsection{OpenVPN}

Jeden z najpopularniejszych protokołów realizujący model scentralizowany ze wsparciem dla TCP i UDP. Często stanowi punkt odniesienia w pracach porównawczych. Udostępniany jest na licencji GPLv2.\\

Autoryzacja opiera się o PKI z obustronną weryfikacją. Certificate Authority podpisuje zarówno certyfikat serwera jak i każdego klienta. Odebranie dostępu odbywa się przez wpisanie na Certificate Revocation Lists.\\

Wspierane są również dodatkowe metody weryfikacji z wykorzystaniem loginu i hasła takie jak: LDAP, RADIUS, PAM, SAML, TOTP/MFA, Active Directory i tym podobne.\\

Dostępny jest oficjalny klient graficzny dla każdego popularnego systemu operacyjnego za wyłączeniem Linux (tylko CLI). Sam serwer jest dostępny w dwóch wariantach: Access Server oraz Community Edition.\\

Access Server jest płatnym produktem nastawionym na klienta korporacyjnego, posiada interfejs webowy posiadających szerokie możliwości konfiguracji oraz wiele narzędzi ułatwiających wdrażanie nowych klientów.\\

Community Edition jest w pełni darmowy, nie posiada interfejsu graficznego, ma jedynie podstawowe funkcjonalności. Jego wdrożenie jest skomplikowane oraz wymaga ręcznego tworzenia wszystkich plików konfiguracyjnych i kryptograficznych.\\

\newpage
\subsection{WireGuard}

Minimalistyczny protokół przewyższający wydajnością OpenVPN w większości publikacji. W przeciwieństwie do OpenVPN konfiguracja jest skrajnie prosta bo oparta o pojedynczy plik konfiguracyjny.\\

Słynie z wyjątkowo krótkiego kodu źródłowego składającego się z około 4 tysięcy lini kodu, kilkudziesięciokrotnie mniej od konkurencyjnych rozwiązań.\\

Autoryzacja bazuje w całości i wyłącznie na kryptografii asymetrycznej. Każdy węzeł posiada statyczną parę kluczy Curve25519.\\

Wspiera jedynie protokół UDP, nie ma żadnej wbudowanej formy automatycznego odkrywania węzłów czy mechanizmów omijania NAT. Wymaga znajomości adresu IP i klucza publicznego wszystkich węzłów, już na poziomie tworzenia konfiguracji.\\

Wszelkie aspekty sieciowe są precyzyjnie definiowane przez poszczególne reguły iptables, nie przez konfigurację samego węzła. Mimo że protokół jest stricte P2P, dowolność i niskopoziomowość konfiguracji umożliwia stworzenie sieci o dowolnej charakterystyce.\\

Surowy charakter protokołu wprowadza wysoki próg wejścia, natomiast staje się dużym atutem po głębszym zapoznaniu się z dokumentacją.\\

\newpage
\subsection{Nebula}

Stosunkowo nowy, bo powstały w 2019 roku protokół wprowadzony przez Slack \cite{OEWV} a obecnie utrzymywany przez Defined Networking \cite{N7BM}. U podstaw buduje sieć typu mesh.\\

Autoryzacja jest oparta o mechanizm typu PKI, który umożliwia obustronną weryfikację podczas tworzenia połączeń pomiędzy węzłami. Para kluczy ED25519 pełni rolę Nebula Certificate Authority która podpisuje certyfikaty poszczególnych węzłów.\\

Każdy certyfikat zawiera informacje takie jak predefiniowane IP wewnątrz sieci, do jakich grup należy węzeł (arbitralne ciągi znaków np. administratorzy) oraz okres ważności certyfikatu (not-before/not-after).\\

Dodatkowa kontrola dostępu (filtrowanie pakietów) na poziomie każdego węzła realizowana jest na podstawie przynależności do grupy.\\

Podobnie jak przy WireGuard, konfiguracja każdego węzła oparta jest o pojedynczy plik konfiguracyjny.\\

W prawie każdej sieci Nebula przynajmniej jeden węzeł ma statyczne IP i pełni rolę lighthouse. Ta funkcjonalność umożliwia innym węzłom automatyczne odkrywanie (peer discovery) oraz asystuje przy przechodzeniu przez NAT (NAT traversal). W przypadku gdy nawiązanie połączenia P2P jest niemożliwe, węzeł lighthouse może opcjonalnie pełnić rolę przekaźnika dla innych węzłów.\\

Brak serwera centralnego, mechanizm automatycznego odkrywania z priorytetyzacją połączeń lokalnych oraz prostota konfiguracji pozytywnie wyróżniają ten protokół pośród innych badanych.